\chapter{Lengua de Señas Hondureña (LESHO)}
\label{cp:lesho}

\insertminitoc
\parindent0pt

El capítulo anterior estableció que las personas sordas conforman una comunidad
lingüística cuya lengua materna es la lengua de señas. En Honduras, esa lengua es el
Lenguaje de Señas Hondureño (\acrshort{lesho}), que constituye a la vez el objeto de reconocimiento
y el medio de comunicación que esta investigación busca apoyar. Este capítulo describe sus
características fundamentales: su historia y reconocimiento oficial, los parámetros que
componen las señas, la distinción entre señas estáticas y dinámicas, el alfabeto
dactilológico y su gramática en relación con el español.

\section{Origen, Evolución y Reconocimiento Oficial}

El Lenguaje de Señas Hondureño (\acrshort{lesho}) es la lengua natural de la comunidad sorda de
Honduras. Su proceso de documentación formal se inició en 1983 con la publicación del
primer diccionario impreso de lengua de señas en el país, denominado ``Las Manos Hablan'',
elaborado en el seno de la Asociación de Sordos de Honduras (\acrshort{ash}). A esta obra le
sucedieron una segunda edición en 2005, titulada ``Comuniquemos Mejor'', y una tercera
en 2019, denominada ``Diccionario de \acrshort{lesho} Tomo I, Primera Edición'', que documenta más
de 900 señas \cite{ash_diccionario_2019}.

El reconocimiento oficial del \acrshort{lesho} se concretó mediante el Decreto Legislativo
No. 321-2013, aprobado por el Congreso Nacional de Honduras, que estableció la Ley de la
Lengua de Señas Hondureña. Esta ley reconoce el \acrshort{lesho} como un derecho en beneficio de
las personas sordas, con discapacidad auditiva y sordociegas que libremente decidan
utilizarla como medio o sistema lingüístico para comunicarse, y obliga al Estado a promover
su uso en los ámbitos educativo, institucional y social \cite{congreso_decreto_2013}.

Williams y Parks documentaron en su encuesta sociolingüística que el \acrshort{lesho} presenta
variaciones regionales, especialmente entre ciudades como Tegucigalpa, San Pedro Sula y
La Ceiba. A pesar de estas variaciones, el estudio determinó que existe un alto grado de
similitud léxica entre las distintas regiones, con coincidencias superiores al 74\,\% en
comparaciones entre diccionarios. La estandarización de la lengua ha sido impulsada
principalmente por la \acrshort{ash} y por los cursos de \acrshort{lesho} ofrecidos en instituciones educativas
del país \cite{williams_encuesta_2015}.

\section{Parámetros Formacionales de las Señas}

Al igual que las lenguas orales se descomponen en unidades mínimas sin significado propio
llamadas fonemas, las lenguas de señas también poseen una estructura fonológica
compuesta por unidades formacionales. Esta noción fue establecida por William Stokoe en
1960, cuyo trabajo demostró que las señas no son gestos holísticos e indivisibles, sino que
se componen de un conjunto finito de elementos sin significado que se combinan y
recombinan para formar el léxico de la lengua. Este principio, conocido como doble
articulación, es una propiedad fundamental que las lenguas de señas comparten con las
lenguas orales \cite{stokoe_structure_1960}.

Stokoe identificó tres parámetros formacionales básicos: la configuración de la mano
(handshape), la ubicación o lugar donde se realiza la seña (location), y el movimiento
(movement). Posteriormente, la investigación lingüística incorporó un cuarto parámetro,
la orientación de la palma (orientation), y reconoció la existencia de componentes no
manuales como las expresiones faciales y los movimientos del cuerpo \cite{sandler_universals_2006}.

Una característica distintiva de las lenguas de señas, señalada por Sandler y Lillo-Martin,
es que estos parámetros se realizan de forma simultánea y no secuencial. Mientras que en
una lengua oral los fonemas se suceden uno tras otro en el tiempo, en una seña la
configuración de la mano, la ubicación, el movimiento y la orientación ocurren al mismo
tiempo. Esta simultaneidad es consecuencia directa de la modalidad visual-gestual de la
lengua y representa una diferencia estructural fundamental respecto a las lenguas orales \cite{sandler_universals_2006}.

El cambio en cualquiera de estos parámetros puede modificar por completo el significado de
una seña, de modo que existen pares mínimos análogos a los de las lenguas orales, donde
dos señas se diferencian únicamente por un parámetro \cite{sandler_universals_2006}. Esta propiedad es de particular
relevancia para el reconocimiento automático, dado que el sistema debe ser capaz de
distinguir variaciones sutiles en la configuración y posición de la mano.

\section{Señas Estáticas y Señas Dinámicas}

Desde la perspectiva del reconocimiento automático, las señas pueden clasificarse según la
presencia o ausencia de movimiento en su ejecución. Las señas estáticas son aquellas que
se definen por una configuración y orientación de la mano fijas en una ubicación
determinada, sin que el movimiento sea un componente esencial de su significado. Un
ejemplo característico de señas estáticas es el alfabeto dactilológico, en el cual cada letra
corresponde a una configuración manual específica que se mantiene durante su
articulación.

Las señas dinámicas, por el contrario, incorporan el movimiento como un parámetro
formacional indispensable. En estas señas, la trayectoria, dirección y velocidad del
movimiento de las manos forman parte integral del significado, por lo que no pueden
representarse adecuadamente mediante una única configuración estática. La mayoría del
vocabulario conversacional del \acrshort{lesho} corresponde a señas dinámicas \cite{stokoe_structure_1960}, \cite{sandler_universals_2006}.

Esta distinción tiene implicaciones técnicas directas en el diseño de sistemas de
reconocimiento. Las señas estáticas pueden clasificarse a partir del análisis de un único
fotograma, ya que toda la información relevante está contenida en una sola pose. Las señas
dinámicas, en cambio, requieren el análisis de una secuencia temporal de fotogramas para
capturar la evolución del movimiento, lo que demanda modelos capaces de procesar
información secuencial.

\section{Alfabeto Dactilológico Hondureño}

El alfabeto dactilológico, también conocido como deletreo manual o dactilología, es un
sistema en el cual cada letra del alfabeto escrito se representa mediante una configuración
manual específica. En el \acrshort{lesho}, el alfabeto dactilológico permite deletrear palabras que
no cuentan con una seña propia, como nombres propios, términos técnicos o palabras
nuevas que aún no se han incorporado al léxico de la lengua \cite{ash_diccionario_2019}.

El alfabeto dactilológico cumple una función complementaria dentro de la comunicación en
lengua de señas. No sustituye al vocabulario de señas, sino que actúa como recurso de
apoyo cuando no existe una seña establecida para un concepto determinado. Williams y
Parks señalaron que la comunidad sorda hondureña recurre al deletreo manual
principalmente para nombres propios y para términos sin seña conocida, mientras que la
comunicación cotidiana se realiza fundamentalmente mediante señas léxicas \cite{williams_encuesta_2015}.

Desde el punto de vista del reconocimiento automático, el alfabeto dactilológico hondureño
está compuesto mayoritariamente por configuraciones manuales estáticas que se ejecutan
con una sola mano, lo que lo convierte en un punto de partida adecuado para sistemas de
reconocimiento basados en la clasificación de poses. Cabe señalar que algunas letras del
alfabeto, como aquellas que representan dígrafos o caracteres específicos del español,
pueden incorporar un componente de movimiento.

\section{Estructura Gramatical de LESHO frente al Español}

El \acrshort{lesho} es una lengua con una gramática propia e independiente del español, y no
constituye una representación manual de este. Esta independencia gramatical es una
característica compartida por todas las lenguas de señas, las cuales, según han demostrado
Sandler y Lillo-Martin, poseen estructuras morfológicas y sintácticas autónomas que no
derivan de las lenguas orales con las que coexisten \cite{sandler_universals_2006}.

Una de las diferencias más relevantes entre el \acrshort{lesho} y el español radica en la organización
de la información. Mientras que el español es una lengua de modalidad auditiva-vocal que
organiza el discurso de forma predominantemente secuencial, el \acrshort{lesho} utiliza el espacio
tridimensional frente al cuerpo del hablante como un recurso gramatical. La ubicación de
las señas en el espacio, la dirección de los movimientos y el uso de referencias espaciales
cumplen funciones gramaticales que en español se expresan mediante palabras o
flexiones \cite{sandler_universals_2006}.

Otra diferencia significativa es el uso de componentes no manuales con valor gramatical.
Las expresiones faciales, la dirección de la mirada y los movimientos de la cabeza y el
cuerpo no son meros acompañamientos expresivos, sino que pueden modificar el significado
de una oración, marcar preguntas, negaciones o condiciones, y cumplir funciones similares
a las de la entonación en las lenguas orales \cite{sandler_universals_2006}.

Estas diferencias estructurales implican que una traducción directa palabra por palabra
entre el español y el \acrshort{lesho} no produce una comunicación gramaticalmente correcta. El
reconocimiento de señas aisladas, como el que aborda el presente trabajo, constituye un
primer nivel de procesamiento que no contempla la totalidad de la complejidad gramatical
de la lengua, aspecto que se reconoce explícitamente dentro de las limitaciones del estudio.
